% Source: physical pp. 366--378 (printed pp. 343--355), beginning at
% the section heading on p. 366 and ending immediately before section
% 31.7 on p. 378.
% Coverage: Eqs. (31.6.1)--(31.6.79), five unnumbered displays, two
% footnotes, three linked citation markers, one centered typographic
% divider, and no figures or tables.
% Source-visible anomalies on physical p. 374 are preserved: Eq.
% (31.6.63) prints \kappa^3, and the following unnumbered display omits
% K from both logarithm arguments.

\subsection{Supergravity to All Orders}

Although the action derived from the Lagrangian density
\eqref{eq:31.2.15} is invariant to zeroth order in
\(\kappa\equiv\sqrt{8\pi G}\) under the local supersymmetry
transformations constructed in Sections 26.7 and 31.5, it is \emph{not}
invariant to first order in \(\kappa\), because the interaction of matter
with the gravitational supermultiplet introduces terms of order
\(\kappa\) in \(\partial_\mu S^\mu\). To make possible the invariance of
the full action, we will have to add terms of higher order in \(\kappa\)
to the Lagrangian and to the supersymmetry transformation rules for the
components of both the matter and the gravitational supermultiplets. This
can be done by first adding terms of higher order in \(\kappa\) to the
transformation rules, chosen so that these transformations form a closed
algebra along with local Lorentz transformations and general coordinate
transformations, and then adding terms to the action to make it invariant
under all these transformations.

This is a long and tedious process. Here we shall give only the
results~\hyperref[ch31-ref-12]{[12]}, and then turn in the next section
to what has been their most important application. The local
supersymmetry transformation of the vierbein, gravitino field, and
auxiliary fields takes the form:
\begin{equation}
  \delta e^a{}_\mu
  =\kappa\bigl(\bar\alpha\gamma^a\psi_\mu\bigr),
  \tag{31.6.1}\label{eq:31.6.1}
\end{equation}
\begin{equation}
  \delta\psi_\mu
  =(2/\kappa)D_\mu\alpha
  +2i\gamma_5\left(
    b_\mu-\frac{1}{3}\gamma_\mu\gamma_\rho b^\rho\right)\alpha
  +\frac{2}{3}\gamma_\mu(s-i\gamma_5p)\alpha,
  \tag{31.6.2}\label{eq:31.6.2}
\end{equation}
\begin{equation}
  \delta s
  =\frac{1}{4e}\bigl(\bar\alpha\gamma_\mu L^\mu\bigr)
  +\frac{\kappa}{2}
   \bigl(\bar\alpha
   [i\gamma_5b^\nu-s\gamma^\nu-ip\gamma_5\gamma^\nu]\psi_\nu\bigr),
  \tag{31.6.3}\label{eq:31.6.3}
\end{equation}
\begin{equation}
  \delta p
  =-\frac{i}{4e}\bigl(\bar\alpha\gamma_5\gamma_\mu L^\mu\bigr)
  +\frac{\kappa}{2}
   \bigl(\bar\alpha
   [b^\nu+is\gamma_5\gamma^\nu-p\gamma^\nu]\psi_\nu\bigr),
  \tag{31.6.4}\label{eq:31.6.4}
\end{equation}
\begin{multline}
  \delta b_\mu
  =\frac{3i}{4e}
    \left(\bar\alpha\gamma_5
      \left[L_\mu-\frac{1}{3}\gamma_\mu\gamma_\rho L^\rho\right]\right)
  +\frac{\kappa}{2}b_\nu
    \bigl(\bar\alpha\gamma^\nu\psi_\mu\bigr)\\
  +\frac{i\kappa}{2}
    \bigl(\bar\psi_\mu\gamma_5(s-i\gamma_5p)\alpha\bigr)
  -\frac{i\kappa}{4}\epsilon_{\mu\nu\kappa\sigma}b^\nu
    \bigl(\bar\alpha\gamma_5\gamma^\kappa\psi^\sigma\bigr).
  \tag{31.6.5}\label{eq:31.6.5}
\end{multline}
Here \(D_\mu\) is again the covariant derivative given by
\eqref{eq:31.5.15} and \eqref{eq:31.5.16}:
\begin{equation}
  D_\mu\equiv
  \partial_\mu+\frac{1}{8}[\gamma_a,\gamma_b]\omega_\mu^{ab},
  \tag{31.6.6}\label{eq:31.6.6}
\end{equation}
but now with the spin connection including terms bilinear in the
gravitino field
\begin{multline}
  \omega_\mu^{ab}
  =e^a{}_\lambda e^b{}_{\nu;\mu}g^{\lambda\nu}
  +\frac{\kappa^2}{4}\left[
    e^b{}_\nu\bigl(\bar\psi_\mu\gamma^a\psi^\nu\bigr)
    +e^a{}_\nu e^b{}_\rho
      \bigl(\bar\psi^\nu\gamma_\mu\psi^\rho\bigr)\right.\\
  \left.
    {}-e^a{}_\nu\bigl(\bar\psi_\mu\gamma^b\psi^\nu\bigr)
  \right].
  \tag{31.6.7}\label{eq:31.6.7}
\end{multline}
Also, \(L^\mu\) is the covariant version of the Rarita--Schwinger
operator \eqref{eq:31.2.8}
\begin{equation}
  L^\mu=i\gamma_5\gamma_\nu D_\rho\psi_\sigma
    \epsilon^{\mu\nu\rho\sigma},
  \tag{31.6.8}\label{eq:31.6.8}
\end{equation}
\(\gamma_\mu\) is defined in terms of the usual Dirac matrices
\(\gamma_a\) as
\begin{equation}
  \gamma_\mu=e^a{}_\mu\gamma_a,
  \tag{31.6.9}\label{eq:31.6.9}
\end{equation}
and \(e\) is here the determinant of the vierbein
\begin{equation}
  e\equiv\sqrt{\operatorname{Det}g}\,.
  \tag{31.6.10}\label{eq:31.6.10}
\end{equation}

It is easy to see that these transformation rules reduce in the weak
field limit to Eqs. \eqref{eq:31.5.2}, \eqref{eq:31.5.19}, and
\eqref{eq:31.5.4}--\eqref{eq:31.5.6}.

The action for pure supergravity that is invariant under these
transformations, and that reduces to the form \eqref{eq:31.2.9} in the
weak field limit, is
\begin{equation}
  I_{\mathrm{SUGRA}}
  =\int d^4x\left[
    -\frac{e}{2\kappa^2}R
    -\frac{1}{2}\bigl(\bar\psi_\mu L^\mu\bigr)
    -\frac{4e}{3}\bigl(s^2+p^2-b_\mu b^\mu\bigr)
  \right],
  \tag{31.6.11}\label{eq:31.6.11}
\end{equation}
where \(R\) is the curvature scalar, calculated using the spin
connection \eqref{eq:31.6.7}. In the absence of matter, the action would
be stationary for \(s=p=b_\mu=0\), leaving us with the simpler action
\[
  I_{\mathrm{SUGRA}}
  =\int d^4x\left[
    -\frac{e}{2\kappa^2}R
    -\frac{1}{2}\bigl(\bar\psi_\mu L^\mu\bigr)
  \right].
\]

The transformation of matter fields is also now more complicated. For
the components of a general scalar supermultiplet, these are
\begin{equation}
  \delta C=i\bigl(\bar\alpha\gamma_5\omega\bigr),
  \tag{31.6.12}\label{eq:31.6.12}
\end{equation}
\begin{equation}
  \delta\omega
  =[-i\gamma_5\sl{\mathcal D}C-M+i\gamma_5N+\sl{V}]\alpha,
  \tag{31.6.13}\label{eq:31.6.13}
\end{equation}
\begin{equation}
  \delta M
  =-\bigl(\bar\alpha[\lambda+\sl{\mathcal D}\omega]\bigr)
  +\frac{2\kappa}{3}
    \bigl(\bar\alpha[s-i\gamma_5p+i\gamma_5\sl{b}]\omega\bigr),
  \tag{31.6.14}\label{eq:31.6.14}
\end{equation}
\begin{equation}
  \delta N
  =i\bigl(\bar\alpha\gamma_5[\lambda+\sl{\mathcal D}\omega]\bigr)
  +\frac{2i\kappa}{3}
    \bigl(\bar\alpha
      [s-i\gamma_5p+i\gamma_5\sl{b}]\gamma_5\omega\bigr),
  \tag{31.6.15}\label{eq:31.6.15}
\end{equation}
\begin{equation}
  \delta V_a
  =\bigl(\bar\alpha\gamma_a\lambda\bigr)
   +\bigl(\bar\alpha\mathcal D_a\omega\bigr)
   +\frac{\kappa}{3}
    \bigl(\bar\alpha
      [s-i\gamma_5p+i\gamma_5\sl{b}]\gamma_a\omega\bigr),
  \tag{31.6.16}\label{eq:31.6.16}
\end{equation}
\begin{equation}
  \delta\lambda
  =-\frac{1}{4}[\gamma^a,\gamma^b]\alpha F_{ab}
   +iD\gamma_5\alpha,
  \tag{31.6.17}\label{eq:31.6.17}
\end{equation}
\begin{equation}
  \delta D=i\bigl(\bar\alpha\gamma_5\sl{\mathcal D}\lambda\bigr),
  \tag{31.6.18}\label{eq:31.6.18}
\end{equation}
where the covariant derivatives here are
\begin{equation}
  \mathcal D_aC
  =e_a{}^\mu\left[
    \partial_\mu C-\frac{i\kappa}{2}
    \bigl(\bar\psi_\mu\gamma_5\omega\bigr)
  \right],
  \tag{31.6.19}\label{eq:31.6.19}
\end{equation}
\begin{multline}
  \mathcal D_a\omega
  =e_a{}^\mu\left[
    \partial_\mu\omega
    +\frac{1}{8}\omega_\mu^{cb}[\gamma_c,\gamma_b]\omega
    -i\kappa b_\mu\gamma_5\omega\right.\\
  \left.
    {}-\frac{\kappa}{2}
      (\sl{V}-i\gamma_5\sl{\mathcal D}C-M+i\gamma_5N)\psi_\mu
  \right],
  \tag{31.6.20}\label{eq:31.6.20}
\end{multline}
\begin{multline}
  \mathcal D_a\lambda
  =e_a{}^\mu\left[
    \partial_\mu\lambda
    +\frac{1}{8}\omega_\mu^{cb}[\gamma_c,\gamma_b]\lambda
    +i\kappa b_\mu\gamma_5\lambda
    +\frac{\kappa}{8}[\gamma_b,\gamma_c]\psi_\mu F_{bc}\right.\\
  \left.
    {}-\frac{i\kappa}{2}\gamma_5D\psi_\mu
  \right],
  \tag{31.6.21}\label{eq:31.6.21}
\end{multline}
\begin{equation}
  F_{ab}
  =e_a{}^\mu e_b{}^\nu\left[
    \partial_\mu V_\nu
    +\frac{\kappa}{2}\partial_\mu(\bar\psi_\nu\omega)
    -\frac{\kappa}{2}(\bar\psi_\mu\gamma_\nu\lambda)
  \right]-a\leftrightarrow b,
  \tag{31.6.22}\label{eq:31.6.22}
\end{equation}
with \(V_\mu\equiv e^a{}_\mu V_a\). The rules for multiplying general
scalar multiplets are the same as Eqs. \eqref{eq:26.2.19}--\eqref{eq:26.2.25},
except that \(\partial_\mu\) is everywhere replaced with \(\mathcal D_a\),
so that the components of a supermultiplet \(S=S_1S_2\) are
\begin{equation}
  C=C_1C_2,
  \tag{31.6.23}\label{eq:31.6.23}
\end{equation}
\begin{equation}
  \omega=C_1\omega_2+C_2\omega_1,
  \tag{31.6.24}\label{eq:31.6.24}
\end{equation}
\begin{equation}
  M=C_1M_2+C_2M_1+\frac{i}{2}
    \bigl(\bar\omega_1\gamma_5\omega_2\bigr),
  \tag{31.6.25}\label{eq:31.6.25}
\end{equation}
\begin{equation}
  N=C_1N_2+C_2N_1-\frac{1}{2}
    \bigl(\bar\omega_1\omega_2\bigr),
  \tag{31.6.26}\label{eq:31.6.26}
\end{equation}
\begin{equation}
  V^a=C_1V_2^a+C_2V_1^a-\frac{i}{2}
    \bigl(\bar\omega_1\gamma_5\gamma^a\omega_2\bigr),
  \tag{31.6.27}\label{eq:31.6.27}
\end{equation}
\begin{multline}
  \lambda
  =C_1\lambda_2+C_2\lambda_1
  -\frac{1}{2}\gamma^a\omega_1\mathcal D_aC_2
  -\frac{1}{2}\gamma^a\omega_2\mathcal D_aC_1
  +\frac{i}{2}\sl{V}_1\gamma_5\omega_2
  +\frac{i}{2}\sl{V}_2\gamma_5\omega_1\\
  +\frac{1}{2}(N_1-i\gamma_5M_1)\omega_2
  +\frac{1}{2}(N_2-i\gamma_5M_2)\omega_1,
  \tag{31.6.28}\label{eq:31.6.28}
\end{multline}
\begin{multline}
  D
  =-\mathcal D_aC_1\mathcal D^aC_2
   +C_1D_2+C_2D_1+M_1M_2+N_1N_2\\
  -\bigl(\bar\omega_1[\lambda_2+\tfrac{1}{2}
      \sl{\mathcal D}\omega_2]\bigr)
  -\bigl(\bar\omega_2[\lambda_1+\tfrac{1}{2}
      \sl{\mathcal D}\omega_1]\bigr)
  -V_{1a}V_2^a.
  \tag{31.6.29}\label{eq:31.6.29}
\end{multline}

Just as in flat space, there are supersymmetric constraints that can be
imposed on these general supermultiplets. One such constraint is
reality. Inspection of Eqs. \eqref{eq:31.6.12}--\eqref{eq:31.6.22}
together with Eqs. \eqref{eq:26.A.20}--\eqref{eq:26.A.21} shows that if
\(C,M,N,V_a,D,\omega,\) and \(\lambda\) form a supermultiplet \(S\),
then \(C^*,M^*,N^*,V_a^*,D^*,\beta\epsilon\gamma_5\omega^*,\) and
\(\beta\epsilon\gamma_5\lambda^*\) are also the components of a
supermultiplet, called \(S^*\). In particular, a \emph{real}
supermultiplet is one with \(S=S^*\), so that \(C,M,N,V_a,\) and \(D\)
are real and \(\omega\) and \(\lambda\) are Majorana spinors.

We can also impose a supersymmetric condition of chirality. Suppose we
set
\begin{equation}
  \lambda=0,\qquad D=0,\qquad
  V_\nu+\frac{1}{2}\kappa(\bar\psi_\nu\omega)=\partial_\nu Z,
  \tag{31.6.30}\label{eq:31.6.30}
\end{equation}
for some field \(Z\). Then Eq. \eqref{eq:31.6.22} gives \(F_{ab}=0\), so
Eq. \eqref{eq:31.6.17} shows that \(\delta\lambda=0\), while
Eq. \eqref{eq:31.6.21} gives \(\mathcal D_a\lambda=0\), so
Eq. \eqref{eq:31.6.18} shows that \(\delta D=0\). The conditions
\(\lambda=D=0\) are therefore preserved by local supersymmetry
transformations. With a little more work one can show that
\begin{equation}
  \delta\left[
    V_\nu+\frac{1}{2}\kappa(\bar\psi_\nu\omega)\right]
  =\partial_\nu(\bar\alpha\omega),
  \tag{31.6.31}\label{eq:31.6.31}
\end{equation}
so the remaining condition, that
\(V_\nu+\frac{1}{2}\kappa(\bar\psi_\nu\omega)\) is a spacetime gradient,
is also supersymmetric. A supermultiplet of component fields satisfying
Eq. \eqref{eq:31.6.30} is said to be \emph{chiral}. Just as in the case
of global supersymmetry, the components of a chiral supermultiplet are
conventionally renamed as \(A,B,\psi,F,\) and \(G\), defined by
\begin{equation}
  C=A,\qquad
  \omega=-i\gamma_5\psi,\qquad
  M=G,\qquad
  N=-F,\qquad
  Z=B.
  \tag{31.6.32}\label{eq:31.6.32}
\end{equation}

A chiral supermultiplet is real if \(A,B,F,\) and \(G\) are real and
\(\psi\) is a Majorana spinor. Such a real chiral supermultiplet can be
written as the sum of a left-chiral supermultiplet \(\Phi\), with
components conventionally defined by
\begin{equation}
  \phi\equiv\frac{A+iB}{\sqrt{2}},\qquad
  \psi_L\equiv\left(\frac{1+\gamma_5}{2}\right)\psi,\qquad
  \mathcal F\equiv\frac{F-iG}{\sqrt{2}},
  \tag{31.6.33}\label{eq:31.6.33}
\end{equation}
and its complex conjugate, a right-chiral supermultiplet. From these
definitions and Eqs. \eqref{eq:31.6.12}--\eqref{eq:31.6.15} and
\eqref{eq:31.6.31} we see that the components of a left-chiral
supermultiplet have the transformation properties
\begin{equation}
  \delta\phi=\sqrt{2}\bigl(\bar\alpha\psi_L\bigr),
  \tag{31.6.34}\label{eq:31.6.34}
\end{equation}
\begin{equation}
  \delta\psi_L
  =\sqrt{2}(\sl{\partial}\phi)\alpha_R
   -\kappa\gamma^\mu(\bar\psi_\mu\psi_L)\alpha_R
   +\sqrt{2}\mathcal F\alpha_L,
  \tag{31.6.35}\label{eq:31.6.35}
\end{equation}
\begin{equation}
  \delta\mathcal F
  =\sqrt{2}\bigl(\bar\alpha\sl{\mathcal D}\psi_L\bigr)
   -\frac{2\kappa}{3}
    \bigl(\bar\alpha[s-ip-i\sl{b}]\psi_L\bigr),
  \tag{31.6.36}\label{eq:31.6.36}
\end{equation}
with \(\mathcal D_a\psi\) given by Eqs. \eqref{eq:31.6.20} and
\eqref{eq:31.6.32}. The rules for multiplying left-chiral
supermultiplets are the same as the corresponding rules
\eqref{eq:26.3.27}--\eqref{eq:26.3.29} in the case of global
supersymmetry: The product of left-chiral supermultiplets \(\Phi_1\)
and \(\Phi_2\) is a left-chiral supermultiplet denoted
\(\Phi_1\Phi_2\), with components
\begin{equation}
  \phi=\phi_1\phi_2,
  \tag{31.6.37}\label{eq:31.6.37}
\end{equation}
\begin{equation}
  \psi_L=\phi_1\psi_{2L}+\phi_2\psi_{1L},
  \tag{31.6.38}\label{eq:31.6.38}
\end{equation}
\begin{equation}
  \mathcal F
  =\phi_1\mathcal F_2+\phi_2\mathcal F_1
   -\bigl(\psi_{1L}^{\mathrm T}\epsilon\psi_{2L}\bigr).
  \tag{31.6.39}\label{eq:31.6.39}
\end{equation}

Now we must consider how to construct actions that are invariant under
local supersymmetry transformations as well as general coordinate
transformations and local Lorentz transformations.
Eqs. \eqref{eq:31.6.18} and \eqref{eq:31.6.21} show that the change
under a supersymmetry transformation of the \(D\)-component of a
general supermultiplet \(S\) is no longer simply a spacetime derivative,
so the integral of such a \(D\)-component is not suitable as a term in
the action. Instead, from an arbitrary superfield \(S\) we can form a
density whose integral \emph{is} supersymmetric
\begin{multline}
  [S]_D\equiv
  e\left[
    D^S-\frac{i\kappa}{2}
      \bigl(\bar\psi^\mu\gamma_\mu\gamma_5\lambda^S\bigr)
    +\frac{4\kappa}{3}
      [-sN^S+pM^S-b^\mu V_\mu^S]
    -\frac{i\kappa}{3}
      \bigl(\bar\omega^S\gamma_5\sl{L}\bigr)\right.\\
  \left.
    {}-\frac{\kappa^2}{4}\epsilon^{\mu\rho\sigma\tau}V_\sigma^S
      \bigl(\bar\psi_\rho\gamma_\tau\psi_\mu\bigr)
    -\frac{\kappa^2}{8}\epsilon^{\mu\rho\sigma\tau}
      \bigl(\bar\omega^S\psi_\sigma\bigr)
      \bigl(\bar\psi_\rho\gamma_\tau\psi_\mu\bigr)
  \right]
  -\frac{2\kappa^2}{3}C^S\mathcal L_{\mathrm{SUGRA}},
  \tag{31.6.40}\label{eq:31.6.40}
\end{multline}
where \(\mathcal L_{\mathrm{SUGRA}}\) is the supergravity Lagrangian
density in Eq. \eqref{eq:31.6.11}:
\begin{equation}
  \mathcal L_{\mathrm{SUGRA}}
  =-\frac{e}{2\kappa^2}R
   -\frac{1}{2}\bigl(\bar\psi_\mu L^\mu\bigr)
   -\frac{4e}{3}\bigl(s^2+p^2-b_\mu b^\mu\bigr).
  \tag{31.6.41}\label{eq:31.6.41}
\end{equation}
Likewise, Eqs. \eqref{eq:31.6.36} and \eqref{eq:31.6.20} show that the
change under a supersymmetry transformation of the
\(\mathcal F\)-component of a left-chiral superfield \(X\) is not a
spacetime derivative, so terms must be added to this
\(\mathcal F\)-term to form a density whose integral is supersymmetric:
\begin{multline}
  [X]_{\mathcal F}
  =e\left[
    \mathcal F^X
    +\frac{\kappa}{\sqrt{2}}
      \bigl(\bar\psi_{\mu R}\gamma^\mu\psi_L^X\bigr)
    +\frac{\kappa^2}{4}
      \bigl(\bar\psi_{\mu R}[\gamma^\mu,\gamma^\nu]\psi_{\nu R}\bigr)
      \phi^X\right.\\
  \left.
    {}+2\kappa(s-ip)\phi^X
  \right].
  \tag{31.6.42}\label{eq:31.6.42}
\end{multline}

There are some supermultiplets of special physical interest. One is the
real non-chiral supermultiplet \(\mathrm I\) whose only non-zero
component is \(C=1\). According to Eq. \eqref{eq:31.6.40}, this
supermultiplet has
\begin{equation}
  [\mathrm I]_D=-\frac{2\kappa^2}{3}\mathcal L_{\mathrm{SUGRA}},
  \tag{31.6.43}\label{eq:31.6.43}
\end{equation}
so this gives nothing new.

A more interesting example is provided by a \emph{left-chiral}
supermultiplet \(\mathrm I\) whose only non-zero component is \(\phi=1\).
According to Eq. \eqref{eq:31.6.42}, this supermultiplet has
\begin{equation}
  \operatorname{Re}[\mathrm I]_{\mathcal F}
  =e\left[
    \frac{\kappa^2}{4}
      \bigl(\bar\psi_{\mu L}[\gamma^\mu,\gamma^\nu]\psi_\nu\bigr)
    +4\kappa s
  \right].
  \tag{31.6.44}\label{eq:31.6.44}
\end{equation}
If a term \(c\int d^4x\,\operatorname{Re}[\mathrm I]_{\mathcal F}\)
appears in the Lagrangian density with a real coefficient \(c\), then
the sum of this term and the supergravity action \eqref{eq:31.6.11} is
stationary with respect to variation of auxiliary fields at
\begin{equation}
  s=3c\kappa/2,\qquad p=b^\mu=0,
  \tag{31.6.45}\label{eq:31.6.45}
\end{equation}
so that after eliminating the auxiliary fields, the action will contain
a cosmological constant term
\begin{equation}
  3\kappa^2c^2\int d^4x\,e,
  \tag{31.6.46}\label{eq:31.6.46}
\end{equation}
corresponding to a vacuum energy density \(-3\kappa^2c^2\). Such a term
is needed to keep the vacuum state Lorentz-invariant when supersymmetry
is spontaneously broken; to cancel the positive vacuum energy density
\(F^2/2\) associated with supersymmetry breaking, we must take
\begin{equation}
  3\kappa^2c^2=F^2/2.
  \tag{31.6.47}\label{eq:31.6.47}
\end{equation}

Looking back at Eq. \eqref{eq:31.6.44} and comparing with
Eq. \eqref{eq:31.3.11}, we see that this gives a gravitino
mass~\hyperref[ch31-ref-13]{[13]}
\begin{equation}
  m_g=c\kappa^2
  =\frac{F\kappa}{\sqrt{6}}
  =\sqrt{\frac{4\pi G F^2}{3}},
  \tag{31.6.48}\label{eq:31.6.48}
\end{equation}
in agreement with our previous result \eqref{eq:31.3.17}.

As an illustration of the use of these formulas that will also provide
results needed in the next section, let us calculate the bosonic part of
the Lagrangian density\footnote{Often the term
\(\mathcal L_{\mathrm{SUGRA}}\) is omitted, with the supergravity
Lagrangian introduced instead by including a constant term
\(-3/\kappa^2\) in \(K(\phi,\phi^*)\). We will not follow this practice;
here, with a conventional normalization of the scalar fields, the term
in \(K(\phi,\phi^*)\) of leading order in \(\kappa\) is
\(\sum_n|\phi_n|^2\).}
\begin{equation}
  \mathcal L
  =\mathcal L_{\mathrm{SUGRA}}
   +\frac{1}{2}[K(\Phi,\Phi^*)]_D
   +2\operatorname{Re}[f(\Phi)]_{\mathcal F},
  \tag{31.6.49}\label{eq:31.6.49}
\end{equation}
where \(\mathcal L_{\mathrm{SUGRA}}\) is here the bosonic part of the
supergravity Lagrangian density \eqref{eq:31.6.41},
\(K(\Phi,\Phi^*)\) is a real function of a set of left-chiral
supermultiplets \(\Phi_n\) and their complex conjugates, and \(f(\Phi)\)
is a function of the \(\Phi_n\) alone. The purely bosonic terms in the
multiplication rules \eqref{eq:31.6.23},
\eqref{eq:31.6.25}--\eqref{eq:31.6.27}, and \eqref{eq:31.6.29} are the
same as in global supersymmetry, so we can use either these rules or the
superspace formalism of Chapter 26 to calculate that the supermultiplet
\(K(\Phi,\Phi^*)\) has the bosonic components
\begin{equation}
  C^K=K(\phi,\phi^*),
  \tag{31.6.50}\label{eq:31.6.50}
\end{equation}
\begin{equation}
  M^K=-2\operatorname{Im}\sum_n\left(
    \frac{\partial K(\phi,\phi^*)}{\partial\phi_n}\mathcal F_n
  \right)+\cdots,
  \tag{31.6.51}\label{eq:31.6.51}
\end{equation}
\begin{equation}
  N^K=-2\operatorname{Re}\sum_n\left(
    \frac{\partial K(\phi,\phi^*)}{\partial\phi_n}\mathcal F_n
  \right)+\cdots,
  \tag{31.6.52}\label{eq:31.6.52}
\end{equation}
\begin{equation}
  V_\mu^K=2\operatorname{Im}\sum_n\left(
    \frac{\partial K(\phi,\phi^*)}{\partial\phi_n}\partial_\mu\phi_n
  \right)+\cdots,
  \tag{31.6.53}\label{eq:31.6.53}
\end{equation}
\begin{equation}
  D^K
  =2\sum_{nm}
    \frac{\partial^2K(\phi,\phi^*)}
         {\partial\phi_n\partial\phi_m^*}
    \left(
      -g^{\mu\nu}\partial_\mu\phi_n\partial_\nu\phi_m^*
      +\mathcal F_n\mathcal F_m^*
    \right)+\cdots,
  \tag{31.6.54}\label{eq:31.6.54}
\end{equation}
where the dots indicate terms involving fermion fields. Also, the rules
\eqref{eq:31.6.37}--\eqref{eq:31.6.39} for multiplying left-chiral
supermultiplets are the same as those in global supersymmetry, so we can
use either these rules or the superfield formalism of Chapter 26 to
calculate that the left-chiral supermultiplet \(f(\Phi)\) has the
bosonic components
\begin{equation}
  \phi^f=f(\phi),
  \tag{31.6.55}\label{eq:31.6.55}
\end{equation}
\begin{equation}
  \mathcal F^f
  =\sum_n\frac{\partial f(\phi)}{\partial\phi_n}\mathcal F_n+\cdots,
  \tag{31.6.56}\label{eq:31.6.56}
\end{equation}
with dots again indicating terms involving fermions. Inserting these
results into Eqs. \eqref{eq:31.6.40} and \eqref{eq:31.6.42} then gives
the bosonic terms of the Lagrangian density \eqref{eq:31.6.49} as
\begin{multline}
  \mathcal L_{\mathrm{bosonic}}
  =\left[
    -\frac{e}{2\kappa^2}R
    -\frac{4e}{3}(s^2+p^2-b_\mu b^\mu)
  \right]
  \left[1-\frac{\kappa^2}{3}K(\phi,\phi^*)\right]\\
  -e\sum_{nm}
    \frac{\partial^2K(\phi,\phi^*)}
         {\partial\phi_n\partial\phi_m^*}
    \left(
      g^{\mu\nu}\partial_\mu\phi_n\partial_\nu\phi_m^*
      -\mathcal F_n\mathcal F_m^*
    \right)\\
  +\frac{4\kappa e}{3}\operatorname{Re}\sum_n
    \frac{\partial K(\phi,\phi^*)}{\partial\phi_n}
    \left[
      \mathcal F_n(s+ip)+ib^\mu\partial_\mu\phi_n
    \right]\\
  +2e\operatorname{Re}\left(
    \sum_n\frac{\partial f(\phi)}{\partial\phi_n}\mathcal F_n
    +2\kappa(s-ip)f(\phi)
  \right).
  \tag{31.6.57}\label{eq:31.6.57}
\end{multline}

Now we must eliminate the auxiliary fields by setting them at values
where the Lagrangian density is stationary.\footnote{Though this
procedure is generally followed, it is not strictly correct, for even
though the Lagrangian density is quadratic in the auxiliary fields, the
coefficients of the terms of second order in the auxiliary fields are
not field-independent. In consequence, in doing the path integration
over auxiliary fields we encounter determinants of the coefficients of
the quadratic terms, which are equivalent to adding terms to the
Lagrangian proportional to
\(\delta^4(0)=(2\pi)^{-4}\int d^4k\,1\). Such terms can be eliminated
by using dimensional regularization, for which \(\int d^4k\,1=0\).}
This gives the auxiliary fields
\begin{multline}
  \mathcal F_n
  =\frac{\kappa^2}{3N}\sum_m
    (\mathcal G^{-1})_{mn}
    \frac{\partial K}{\partial\phi_m^*}
    \left[
      -\sum_{k\ell}(\mathcal G^{-1})_{k\ell}
       \frac{\partial K}{\partial\phi_\ell}
       \left(\frac{\partial f}{\partial\phi_k}\right)^*
      +3f^*
    \right]\\
  -\sum_m(\mathcal G^{-1})_{mn}
    \left(\frac{\partial f}{\partial\phi_m}\right)^*,
  \tag{31.6.58}\label{eq:31.6.58}
\end{multline}
\begin{equation}
  s-ip
  =\frac{\kappa}{2N}\left[
    -\sum_{\ell k}(\mathcal G^{-1})_{k\ell}
      \frac{\partial K}{\partial\phi_\ell}
      \left(\frac{\partial f}{\partial\phi_k}\right)^*
    +3f^*
  \right],
  \tag{31.6.59}\label{eq:31.6.59}
\end{equation}
\begin{equation}
  b_\mu
  =\frac{\kappa}{2(1-\kappa^2K/3)}
    \operatorname{Im}\left(
      \sum_n\frac{\partial K}{\partial\phi_n}\partial_\mu\phi_n
    \right),
  \tag{31.6.60}\label{eq:31.6.60}
\end{equation}
where
\[
  N\equiv
  1-\frac{\kappa^2}{3}K
  +\frac{\kappa^2}{3}\sum_{k\ell}
    (\mathcal G^{-1})_{k\ell}
    \frac{\partial K}{\partial\phi_\ell}
    \frac{\partial K}{\partial\phi_k^*}
\]
and \(\mathcal G(\phi,\phi^*)\) is the Kahler metric
\begin{equation}
  \mathcal G_{nm}(\phi,\phi^*)
  \equiv
  \frac{\partial^2K(\phi,\phi^*)}
       {\partial\phi_n\partial\phi_m^*}.
  \tag{31.6.61}\label{eq:31.6.61}
\end{equation}

Using this in Eq. \eqref{eq:31.6.57} gives the bosonic Lagrangian
\begin{multline}
  \mathcal L_{\mathrm{bosonic}}
  =-\frac{e}{2\kappa^2}R
    \left[1-\frac{\kappa^2}{3}K\right]
  -e\sum_{nm}\mathcal G_{nm}g^{\mu\nu}
    \partial_\mu\phi_n\partial_\nu\phi_m^*\\
  +\frac{e\kappa^2}{3N}
    \left|
      \sum_{mn}(\mathcal G^{-1})_{nm}
      \frac{\partial f}{\partial\phi_m}
      \frac{\partial K}{\partial\phi_n^*}
      -3f
    \right|^2\\
  -\frac{e\kappa^2}{3(1-\kappa^2K/3)}
    \operatorname{Im}\left[
      \sum_n\frac{\partial K}{\partial\phi_n}\partial_\mu\phi_n
    \right]
    \operatorname{Im}\left[
      \sum_n\frac{\partial K}{\partial\phi_n}\partial_\nu\phi_n
    \right]g^{\mu\nu}\\
  -e\sum_{mn}(\mathcal G^{-1})_{nm}
    \frac{\partial f}{\partial\phi_m}
    \left(\frac{\partial f}{\partial\phi_n}\right)^*.
  \tag{31.6.62}\label{eq:31.6.62}
\end{multline}
As already noted in the weak field case in Section 31.2, the Lagrangian
density \eqref{eq:31.6.62} has the uncomfortable feature that the
Einstein--Hilbert term \(-eR/2\kappa^2\) is multiplied with a factor
\((1-\kappa^2K(\phi,\phi^*)/3)\), so that the effective gravitational
constant varies from point to point in spacetime. To remedy this, we
perform a Weyl transformation, defining a new metric by
\begin{equation}
  \widetilde g_{\mu\nu}
  =(1-\kappa^3K/3)g_{\mu\nu}.
  \tag{31.6.63}\label{eq:31.6.63}
\end{equation}
The Einstein Lagrangian density is given in terms of the new metric by
\[
  eg^{\mu\nu}R_{\mu\nu}
  =\left(1-\frac{\kappa^2}{3}K\right)^{-1}
    \widetilde e\,\widetilde g^{\mu\nu}
    \left[
      \widetilde R_{\mu\nu}
      +\frac{3}{2}\partial_\mu
        \ln\left(1-\frac{\kappa^2}{3}\right)
        \partial_\nu
        \ln\left(1-\frac{\kappa^2}{3}\right)
    \right],
\]
where \(\widetilde R_{\mu\nu}\) is the curvature tensor calculated using
a metric \(\widetilde g_{\mu\nu}\) in place of \(g_{\mu\nu}\), and
\(\widetilde e\equiv\sqrt{\operatorname{Det}\widetilde g}\). A
straightforward calculation then gives the bosonic Lagrangian
\eqref{eq:31.6.62} as
\begin{multline}
  \widetilde{\mathcal L}_{\mathrm{bosonic}}
  =-\frac{\widetilde e}{2\kappa^2}
    \widetilde g^{\mu\nu}\widetilde R_{\mu\nu}\\
  -\widetilde e\,\widetilde g^{\mu\nu}
    \sum_{nm}\partial_\mu\phi_n\partial_\nu\phi_m^*
    \left[
      \left(1-\frac{\kappa^2}{3}K\right)^{-1}
      \frac{\partial^2K}{\partial\phi_n\partial\phi_m^*}
      +\frac{\kappa^2}{3}
      \left(1-\frac{\kappa^2}{3}K\right)^{-2}
      \frac{\partial K}{\partial\phi_n}
      \frac{\partial K}{\partial\phi_m^*}
    \right]\\
  +\frac{\widetilde e\kappa^2}{3N}
    \left(1-\frac{\kappa^2}{3}K\right)^{-2}
    \left|
      \sum_{mn}(\mathcal G^{-1})_{nm}
      \frac{\partial f}{\partial\phi_m}
      \frac{\partial K}{\partial\phi_n^*}
      -3f
    \right|^2\\
  -\widetilde e
    \left(1-\frac{\kappa^2}{3}K\right)^{-2}
    \sum_{mn}(\mathcal G^{-1})_{nm}
    \frac{\partial f}{\partial\phi_m}
    \left(\frac{\partial f}{\partial\phi_n}\right)^*.
  \tag{31.6.64}\label{eq:31.6.64}
\end{multline}
The Weyl transformation has not only removed the factor
\((1-\kappa^2K/3)\) from the Einstein--Hilbert term; it has also
eliminated terms proportional to
\(\partial_\mu\phi_n\partial_\nu\phi_m\) and
\(\partial_\mu\phi_n^*\partial_\nu\phi_m^*\).

This result may be further simplified by introducing a \emph{modified
Kahler potential} \(d(\phi,\phi^*)\) in place of \(K(\phi,\phi^*)\),
which we define by
\begin{equation}
  1-\frac{\kappa^2}{3}K
  \equiv\exp\left(-\frac{\kappa^2d}{3}\right).
  \tag{31.6.65}\label{eq:31.6.65}
\end{equation}
We also introduce a new metric on the scalar field space
\begin{equation}
  g_{nm}\equiv
  \frac{\partial^2d}{\partial\phi_n\partial\phi_m^*}.
  \tag{31.6.66}\label{eq:31.6.66}
\end{equation}
The reciprocals of the new and old metrics are related by
\[
  \mathcal G^{-1}_{\ell k}
  =\exp(\kappa^2d/3)\left[
    g^{-1}_{\ell k}
    +\frac{\kappa^2}{3}
     \frac{
       \displaystyle\sum_{mn}
       g^{-1}_{\ell n}g^{-1}_{mk}
       (\partial d/\partial\phi_n)
       (\partial d/\partial\phi_m^*)
     }{
       \displaystyle
       1-(\kappa^2/3)\sum_{mn}
       g^{-1}_{mn}
       (\partial d/\partial\phi_n)
       (\partial d/\partial\phi_m^*)
     }
  \right].
\]
The bosonic Lagrangian density \eqref{eq:31.6.64} now takes the simpler
form
\begin{equation}
  \widetilde{\mathcal L}_{\mathrm{bosonic}}
  =-\frac{\widetilde e}{2\kappa^2}
    \widetilde g^{\mu\nu}\widetilde R_{\mu\nu}
   -\widetilde e\,\widetilde g^{\mu\nu}
    \sum_{nm}g_{nm}\partial_\mu\phi_n\partial_\nu\phi_m^*
   -\widetilde eV,
  \tag{31.6.67}\label{eq:31.6.67}
\end{equation}
where \(V(\phi,\phi^*)\) is the potential
\begin{equation}
  V
  =\exp(\kappa^2d)\left[
    \sum_{nm}g^{-1}_{nm}L_mL_n^*
    -3\kappa^2|f|^2
  \right]
  \tag{31.6.68}\label{eq:31.6.68}
\end{equation}
and
\begin{equation}
  L_m\equiv
  \frac{\partial f}{\partial\phi_m}
  +\kappa^2f\frac{\partial d}{\partial\phi_m}.
  \tag{31.6.69}\label{eq:31.6.69}
\end{equation}

The potential \eqref{eq:31.6.68} has an obvious stationary point at
field strengths satisfying the condition \(L_m=0\). However, as we
found in the weak field case, at this point the vacuum energy in general
takes the negative value \(-3\kappa^2|f|^2\). To have the stationary
point \(L_m=0\) give a solution with flat space, it is necessary that
both \(f(\phi)\) and \(\partial f(\phi)/\partial\phi_n\) should vanish at
these field values. Inspection of Eqs. \eqref{eq:31.6.58} and
\eqref{eq:31.6.59} shows that the scalar auxiliary fields
\(\mathcal F_n,s,\) and \(p\) vanish for such field values, so that the
vacuum expectation values of the variations \eqref{eq:31.6.2} and
\eqref{eq:31.6.35} of the gravitino and chiral spinor fields under
global supersymmetry transformations with constant \(\alpha\) vanish.
Thus a vacuum field value for which \(f(\phi)\) and
\(\partial f(\phi)/\partial\phi_n\) all vanish is one for which global
supersymmetry is unbroken in the classical limit. In the next section
we shall consider vacuum configurations in which supersymmetry
\emph{is} broken.

We will not show it here, but the additional terms in the bosonic
Lagrangian required by the inclusion of gauge superfields are
unaffected by gravitation, aside from an over-all determinantal factor
\(\widetilde e\) and metric factors needed to raise and lower indices.
After elimination of auxiliary fields and a Weyl transformation, the
complete bosonic Lagrangian for a theory with gauge as well as chiral
and gravitational superfields is
\begin{multline}
  \mathcal L_{\mathrm{bosonic}}/\widetilde e
  =-\frac{1}{2\kappa^2}\widetilde R^\mu{}_\mu
   -\sum_{nm}g_{nm}D_\mu\phi_nD^\mu\phi_m^*
   -\frac{1}{4}\sum_{AB}\operatorname{Re}f_{AB}
     F_{\mu\nu}^AF^{B\mu\nu}\\
  -\frac{1}{8}\sum_{AB}\operatorname{Im}f_{AB}
     F_{\mu\nu}^AF_{\rho\sigma}^B\epsilon^{\mu\nu\rho\sigma}
  -V.
  \tag{31.6.70}\label{eq:31.6.70}
\end{multline}
Here \(D_\mu,F_{\mu\nu}^A,\) and \(t_A\) denote the gauge-covariant
derivatives, the field-strength tensors, and the representatives of the
gauge generators on the chiral scalar superfields, using the notation
described in Section 15.1; \(f_{AB}\) is an independent holomorphic
function of the \(\phi_n\); all spacetime indices are raised and lowered
with \(\widetilde g_{\mu\nu}\); and the potential \(V\) now takes the
form
\begin{multline}
  V
  =\exp(\kappa^2d)\left[
    \sum_{nm}g^{-1}_{nm}L_mL_n^*
    -3\kappa^2|f|^2
  \right]\\
  +\frac{1}{2}\operatorname{Re}\sum_{AB}f^{-1}_{AB}
    \left(
      \sum_{nm}\frac{\partial d}{\partial\phi_n}
      (t_A)_{nm}\phi_m
    \right)^*
    \left(
      \sum_{k\ell}\frac{\partial d}{\partial\phi_k}
      (t_B)_{k\ell}\phi_\ell
    \right).
  \tag{31.6.71}\label{eq:31.6.71}
\end{multline}
The form \eqref{eq:31.6.71} of the bosonic potential is simple enough to
make it apparent that terms in \(d\) that depend only on \(\phi_n\) or
only on \(\phi_n^*\) may be traded for corrections to the
superpotential. Specifically, if we write
\begin{equation}
  d(\phi,\phi^*)
  =\widetilde d(\phi,\phi^*)+a(\phi)+a(\phi)^*,
  \qquad
  f(\phi)=\widetilde f(\phi)\exp\bigl(-\kappa^2a(\phi)\bigr),
  \tag{31.6.72}\label{eq:31.6.72}
\end{equation}
with \(a(\phi)\) an arbitrary holomorphic function satisfying the gauge
invariance condition
\[
  \sum_{nm}\frac{\partial a(\phi)}{\partial\phi_n}
    (t_A)_{nm}\phi_m=0,
\]
then the potential \eqref{eq:31.6.71} takes the same form in terms of
\(\widetilde d\) and \(\widetilde f\) as it did in terms of \(d\) and
\(f\). With a suitable redefinition of the superpotential, we can then
eliminate any terms in the power series expansion of \(d(\phi,\phi^*)\)
that depend only on \(\phi_n\) or only on \(\phi_n^*\). With this
understanding, the leading term in the power series expansion of
\(d(\phi,\phi^*)\) (now dropping the tilde) is of the form
\(\sum_{nm}d_{nm}\phi_n\phi_m^*\). By a suitable linear transformation
of the superfields, we can then make the matrix \(d_{nm}\) equal to
\(\delta_{nm}\), so that the power series expansion of
\(d(\phi,\phi^*)\) begins
\begin{equation}
  d(\phi,\phi^*)=\sum_n|\phi_n|^2+\cdots,
  \tag{31.6.73}\label{eq:31.6.73}
\end{equation}
and the power series expansion of the metric \eqref{eq:31.6.66} begins
\begin{equation}
  g_{nm}=\delta_{nm}+\cdots.
  \tag{31.6.74}\label{eq:31.6.74}
\end{equation}
Inspection of the second term on the right-hand side of
Eq. \eqref{eq:31.6.67} shows that the scalar fields defined in this way
are canonically normalized.

The fermion terms are much more complicated. Here we will quote only
the terms quadratic in the gaugino fields
\begin{multline}
  \mathcal L_{\mathrm{gaugino}}^{(2)}/\widetilde e
  \equiv
  -\frac{1}{2}\operatorname{Re}\sum_{AB}f_{AB}
    \bigl(\bar\lambda_A\sl{\mathcal D}\lambda_B\bigr)\\
  +\frac{1}{2}\exp(\kappa^2d/2)
    \operatorname{Re}\sum_{mn}\sum_{AB}
    g^{-1}_{nm}L_m
    \left(\frac{\partial f_{AB}}{\partial\phi_n}\right)^*
    \bigl(\bar\lambda_A\lambda_B\bigr),
  \tag{31.6.75}\label{eq:31.6.75}
\end{multline}
with \(L_m\) given by Eq. \eqref{eq:31.6.69}. We see that if the gauge
fields are canonically normalized, then the constant term in the
expansion of the function \(f_{AB}\) in powers of the scalar fields is
\(\delta_{AB}\), and then the gaugino fields \(\lambda_A\) are also
canonically normalized.

\begin{center}
  \(\ast\quad\ast\quad\ast\)
\end{center}

Instead of moving all the holomorphic terms in \(d(\phi,\phi^*)\) and
their complex conjugates into the superpotential, we can use the
transformation \eqref{eq:31.6.72} to make the new superpotential
\(\widetilde f(\phi)\) equal to a constant, which can be chosen to be
equal to unity, by taking
\(a(\phi)=-\kappa^{-2}\ln f(\phi)\). The potential then depends only on
the function
\begin{equation}
  \mathcal D(\phi,\phi^*)
  \equiv d(\phi,\phi^*)+2\kappa^{-2}\operatorname{Re}\ln f(\phi),
  \tag{31.6.76}\label{eq:31.6.76}
\end{equation}
and takes the form
\begin{multline}
  V
  =\exp(\kappa^2\mathcal D)\left[
    \kappa^4\sum_{nm}g^{-1}_{nm}
      \left(\frac{\partial\mathcal D}{\partial\phi_m}\right)
      \left(\frac{\partial\mathcal D}{\partial\phi_n}\right)^*
    -3\kappa^2
  \right]\\
  +\frac{1}{2}\operatorname{Re}\sum_{AB}f^{-1}_{AB}
    \left(
      \sum_{nm}\frac{\partial\mathcal D}{\partial\phi_n}
      (t_A)_{nm}\phi_m
    \right)
    \left(
      \sum_{k\ell}\frac{\partial\mathcal D}{\partial\phi_k}
      (t_B)_{k\ell}\phi_\ell
    \right)^*.
  \tag{31.6.77}\label{eq:31.6.77}
\end{multline}
Also, the metric \eqref{eq:31.6.66} for the scalar fields may be written
\begin{equation}
  g_{nm}
  =\frac{\partial^2\mathcal D}
         {\partial\phi_n\partial\phi_m^*}.
  \tag{31.6.78}\label{eq:31.6.78}
\end{equation}
Although we have not shown it here, the symmetry that allows us to
replace the Kahler potential and superpotential with the single
function \(\mathcal D(\phi,\phi^*)\) also allows us to make this
replacement in the whole Lagrangian, including all terms involving
fermions and gauge fields.

There is an interesting class of `no-scale'
theories~\hyperref[ch31-ref-13a]{[13a]}, in which the potential \(V\)
vanishes for all values of \(\phi_m\). For instance, this is the case
for a single gauge-neutral chiral scalar superfield, with
\begin{equation}
  \mathcal D
  =-3\kappa^{-2}\ln\bigl(h(\phi)+h(\phi)^*\bigr),
  \tag{31.6.79}\label{eq:31.6.79}
\end{equation}
where \(h(\phi)\) is an arbitrary function of \(\phi\). But there is no
known principle that would require \(\mathcal D\) to take this form.
